\section{Non-normal canonical states in double-logarithmic depth}
\label{sec:ldp_nonnormal_trees}

The paragraph ``Tree-RG circuit with measurements'' of
\cite{Malz2024Preparation} gives double-logarithmic depth for short-range
correlated states. Its paragraph ``Long-range MPS using measurements''
prepares coherent sector labels before the blocked isometries. Combining
these constructions on their common virtual support extends the
measurement-assisted depth bound to the canonical non-normal class below.
This is an enhancement of those constructions, not a claim that the source
already proves the all-length statement in this form.

\begin{theorem}[Inequivalence gives the mixed eigenvalue gap]
    \label{thm:ldp_inequivalent_mixed_gap}
    \lean{MPSTensor.mixedMap_eigenvalue_norm_lt_one_of_spectralRadius_lt_one,
        MPSTensor.mixedMap_eigenvalue_norm_lt_one_of_normal_inequivalent}
    \leanok
    Let $A_j,A_k$ be normal left-canonical blocks of positive bond dimensions.
    If their dimensions agree, assume they are not gauge-phase equivalent.
    Every eigenvalue of their mixed transfer map has modulus less than one.
\end{theorem}
\begin{proof}
    \leanok
    Normality implies irreducibility. The mixed-transfer spectral-radius
    theorem gives a strict gap, both for inequivalent square blocks and for
    different bond dimensions. An eigenvalue belongs to the spectrum of
    the mixed map and hence its modulus is bounded by that spectral radius.
\end{proof}

\begin{theorem}[Coherent supported block trees from sparse labels]
    \label{thm:ldp_coherent_supported_blocks}
    \lean{QuantumCircuit.IsPreparedWithMeasurementRoundsInDepth.apply_implementation,
        MPSPreparation.exists_pairFamilyUnitary,
        MPSPreparation.exists_isLocalCircuitOfDepth_registerCfg_of_pairProductState_supported,
        MPSPreparation.pairLayerOp_mulVec_registerCfg_eq,
        MPSPreparation.blockLayerOp_pairLayerOp_mulVec_registerCfg,
        MPSPreparation.exists_isPreparedWithMeasurementRoundsInDepth_sum_blockIsometryState,
        MPSPreparation.exists_supported_block_preparation_constants}
    \leanok
    \uses{thm:ldp_sparse_window_ghz, thm:ldp_supported_tree_rounds,
        lem:ldp_balanced_widths}
    Fix $d\ge2$, a tensor $B$, a set $S$ of virtual pairs, an eventual
    support-injectivity threshold $L$, and a bound on the number of sector
    labels. There are $s\ge\max(2,L)$ and $C$, fixed before the ring and
    amplitudes, with the following property. Suppose
    \begin{align*}
        2^{h+2}s\le\ell_k\le2^{h+1}(8s).
    \end{align*}
    Let the pairs $\omega_j$ be orthonormal and $\sum_j|\alpha_j|^2=1$.
    Assume that, at the actual number $M$ of blocks, each product of the
    cyclic pairs $\omega_j$ is supported on $S^M$. Then
    \begin{align*}
        \sum_j\alpha_j\Bigl(\bigotimes_kV_{\ell_k}\Bigr)
        \bigotimes_k\ket{\omega_j}_{R_kL_{k+1}}
    \end{align*}
    is prepared in quantum depth at most $C(h+1)$.
\end{theorem}
\begin{proof}
    \leanok
    The sparse GHZ construction has depth independent of the block lengths.
    A single bounded-window unitary maps every encoded label $j$ to
    $\omega_j$, preserving all coherent amplitudes. The bond encodings have
    a final zero site. Thus their full physical configurations agree exactly
    with the tree-root windows even when an odd block leaves one terminal
    site outside the even leaf cutoff. The pair layer leaves all tree
    interior scratch sites zero, so the supported compiler applies.
    Compose the actual rounds and add their depths. For a fixed outcome
    history, its scalar is common to all inputs in the zero-scratch
    subspace. The endpoint specifies every physical site.
\end{proof}

\begin{theorem}[All-length compilation of exponential block approximations]
    \label{thm:ldp_supported_loglog_compilation}
    \lean{MPSPreparation.exists_log_log_preparation_of_supported_block_approximation,
        MPSPreparation.natCast_mul_succ_le_mul_log,
        MPSPreparation.exists_two_pow_le}
    \leanok
    \uses{thm:ldp_coherent_supported_blocks, lem:ldp_block_length_threshold_error}
    Under the fixed tensor, support and eventual support-injectivity
    assumptions of Theorem~\ref{thm:ldp_coherent_supported_blocks},
    suppose a normalized target family has coherent supported block
    approximations of error at most $KM e^{-rq}$ for block lengths
    $\ell_k\ge q\ge L$, with $K,r>0$ independent of the amplitudes and
    lengths. Suppose also that every admissible target of length $N\ge L$
    has an exact normalized one-block supported pair. Then a fixed $c$
    gives, for every admissible $N\ge2$ and $0<\varepsilon\le1$, a unit
    approximant of error at most $\varepsilon$ prepared in depth
    \begin{align*}
        T\le c\log\bigl(1+\log(N/\varepsilon)\bigr).
    \end{align*}
    The one-block support assumption is required only for $M=1$.
\end{theorem}
\begin{proof}
    \leanok
    Take $q=\lceil a\log(N/\varepsilon)+b\rceil$, where $a=1/r$ and
    $b\ge\max(\log K,0)/r+4s+1$. If $q\le N$, use $M=\lfloor N/q\rfloor$
    blocks, the last of length $q+N\bmod q$ and all others of length $q$.
    Choose $h$ with $2^{h+2}s\le q<2^{h+1}(4s)$. All block lengths satisfy
    the previous theorem, and $KM e^{-rq}\le KN e^{-rq}\le\varepsilon$.

    If $4s\le N<q$, use a whole-ring supported tree, with $h$ chosen from
    $N$ rather than from $q$. Its root pair gives the target exactly.
    In both branches $2^h\le q\le a\log(N/\varepsilon)+b+1$, so the
    depth $C(h+1)$ has the required bound. For $N<4s$, exact synthesis
    on finitely many possible lengths is absorbed into the constant,
    since $\log(1+\log(N/\varepsilon))\ge\log(1+\log2)>0$.

    % Formula: normalized target = Vhat_N eta_N on the one-block support.
    % Source: supported polar factorization and the exact one-block pair.
    % Boundary signature on each side: no input and one H_N physical output.
    % Internal contraction: a chi-dimensional supported root state into Vhat_N.
    % Vhat_N is the entire bounded-leaf tree, never a free N-site gate.
    \begin{center}
        \tenkzkernel
        \begin{tenkzequation}
            \begin{tenkz}[rows={wire}, cols=1, bonds=none]
                \tn[skin=box, ports={90:physical:$\mathcal H_N$}]{\widehat\phi_N}
            \end{tenkz}
            =
            \begin{tenkz}[rows={wire,wire}, cols=1, bonds=none]
                \tn[at=(1,1), skin=box, name=tree,
                    ports={90:physical:$\mathcal H_N$,270:virtual}]{\widehat V_N}
                \tn[at=(2,1), skin=box, name=root,
                    ports={90:virtual:$\C^\chi$}]{\eta_N}
                \tnwire{root.90}{tree.270}
            \end{tenkz}
        \end{tenkzequation}
    \end{center}
    Here $\widehat\phi_N=\phi_N/\|\phi_N\|$, and $\eta_N$ is the
    normalized supported one-block input, including
    the cyclic ordering of the two virtual legs. The box $\widehat V_N$
    denotes the entire bounded-leaf tree, not one free $N$-site gate.
\end{proof}

\begin{theorem}[Non-normal canonical preparation in double-logarithmic depth]
    \label{thm:ldp_nonnormal_loglog}
    \lean{MPSPreparation.exists_isPreparedWithMeasurementRoundsInDepth_le_log_log_of_mpv_eq_sum,
        MPSPreparation.exists_isPreparedWithMeasurementRoundsInDepth_repeatedBlockSum_le_log_log,
        MPSPreparation.exists_isPreparedWithMeasurementRoundsInDepth_repeatedBlockSum_le_log_log_of_inequivalent}
    \leanok
    \uses{thm:ldp_inequivalent_mixed_gap, thm:ldp_supported_loglog_compilation,
        def:ldp_repeated_block_sum, lem:ldp_common_mixed_gap,
        lem:ldp_overlapping_block_injectivity, thm:ldp_overlapping_unequal_blocks_error,
        lem:ldp_overlapping_one_block_chain}
    Let
    \begin{align*}
        A^i=\bigoplus_j\operatorname{diag}(\mu_{j,1},\ldots,\mu_{j,m_j})
        \otimes A_j^i
    \end{align*}
    be the canonical repeated-block form, with normal inequivalent blocks
    in the left-canonical gauge, and positive definite transfer fixed points
    $\sigma_j$ of trace one. The copy weights are arbitrary nonzero complex
    numbers under the existing convention that zero padding is omitted.
    There is $c$, fixed before $N$ and $\varepsilon$, such that every
    $N\ge2$ with $\phi_N(A)\ne0$ and every $0<\varepsilon\le1$ admit
    a unit vector $\psi$ satisfying
    \begin{align*}
        1-|\langle\psi,\phi_N(A)/\|\phi_N(A)\|\rangle|\le\varepsilon,
        \qquad T\le c\log\bigl(1+\log(N/\varepsilon)\bigr).
    \end{align*}
    The protocol prepares $\psi$ in measurement rounds of quantum depth $T$.
    Finite-block physical ranges may overlap. Individual sums
    $\beta_j(N)=\sum_k\mu_{j,k}^N$ may vanish.
\end{theorem}
\begin{proof}
    \leanok
    Work with the unweighted direct sum $B=\bigoplus_jA_j$ and its
    diagonal sector-pair support. The mixed gap gives eventual support
    injectivity and the amplitude-uniform overlap estimate. Set
    $\alpha_j=\beta_j/(\sum_l|\beta_l|^2)^{1/2}$; nonvanishing of the
    target implies that the coefficient vector is nonzero, without
    requiring each coefficient to be nonzero. The embedded fixed-point
    pairs are orthonormal and their cyclic products have the required
    support. The exact one-block pair supplies the short accuracy branch.
    Apply Theorem~\ref{thm:ldp_supported_loglog_compilation}.

    The intermediate theorem accepts any supplied identity
    $\phi_N(A)=\sum_j\beta_j(N)\phi_N(A_j)$, so the coefficients need not
    be powers for that interface. For the repeated-block form the trace
    identity gives precisely the power sums above.
\end{proof}

Only nearest-neighbour unitary layers contribute to $T$. Measurements,
classical communication and on-site corrections are free, the number of
rounds may grow with the tree height, and no outcome is postselected.
The theorem concerns the supplied canonical block class. Transport from
an arbitrary original tensor, including its periodicity and blocking data,
is a separate statement.
